\documentclass[10pt,a4paper]{article}
\usepackage[T1]{fontenc}
\usepackage[utf8]{inputenc}
\usepackage[ngerman]{babel}
\usepackage[a4paper,margin=14mm]{geometry}
\usepackage{lmodern}
\usepackage[table]{xcolor}
\usepackage{tabularx}
\usepackage{amsmath}
\usepackage{microtype}
\usepackage{tikz}
\usetikzlibrary{arrows.meta,calc}
\newif\ifsolution
\solutionfalse
\pagestyle{empty}
\setlength{\parindent}{0pt}
\setlength{\parskip}{0pt}
\renewcommand{\familydefault}{\sfdefault}
\definecolor{Navy}{HTML}{17324D}
\definecolor{Blue}{HTML}{3B6EA5}
\definecolor{Sky}{HTML}{EAF3FA}
\definecolor{Gold}{HTML}{F3B33D}
\definecolor{Ink}{HTML}{1E2935}
\definecolor{Muted}{HTML}{617080}
\definecolor{Line}{HTML}{D8E1E8}
\color{Ink}
\newcommand{\FZ}{F_{\mathrm Z}}
\newcommand{\FG}{F_{\mathrm G}}
\newcommand{\FS}{F_{\mathrm S}}
\newcommand{\kopf}[1]{%
  \colorbox{Navy}{\parbox{\dimexpr\linewidth-2\fboxsep\relax}{%
    \vspace{1.0mm}\color{white}
    \begin{tabularx}{\linewidth}{@{}Xr@{}}
      {\small\bfseries PHYSIK \textperiodcentered{} KLASSE 11} &
      \colorbox{Gold}{\color{Navy}\small\bfseries\strut\ifsolution LÖSUNGEN\else ABFRAGE\fi}
    \end{tabularx}
    \vspace{0.7mm}
    {\fontsize{16}{19}\selectfont\bfseries Aufgabe 6 -- Kräfte beim Kettenkarussell}\par
    \vspace{0.3mm}{\small #1}\vspace{0.8mm}
  }}\par\vspace{0.4mm}}
\newcommand{\frage}[2]{%
  \vspace{0.2mm}{\color{Blue}\bfseries\large #1 \textperiodcentered{} #2}\par
  \vspace{0.1mm}{\color{Blue}\rule{\linewidth}{0.65pt}}\par\vspace{0.5mm}}
\newcommand{\antwortfeld}[2]{%
  \par\vspace{0.4mm}\noindent
  \pgfmathsetmacro{\aw}{\linewidth/1mm-0.6}%
  \pgfmathsetmacro{\tw}{\aw-1}%
  \begin{tikzpicture}[x=1mm,y=1mm]
    \path[use as bounding box] (0,0) rectangle (\aw,{-9*#1});
    \ifsolution
      \node[anchor=north west,text width=\tw mm,inner sep=0pt,align=left,
        font=\small,text=Blue] at (0,-0.5) {#2};
    \else
      \foreach \i in {1,...,#1}{\draw[Line,line width=.55pt] (0,{-9*\i}) -- (\aw,{-9*\i});}
    \fi
  \end{tikzpicture}\par}
\newcommand{\fuss}[1]{%
  \vfill{\color{Line}\rule{\linewidth}{0.5pt}}\par\vspace{1mm}
  \begin{tabularx}{\linewidth}{@{}Xr@{}}
    {\small\color{Muted}Klasse 11 \textperiodcentered{} Physik \textperiodcentered{} Kreisbewegungen} &
    {\small\color{Muted}\ifsolution Lösungen\else Abfrage\fi{} \textperiodcentered{} Seite #1 von 2}
  \end{tabularx}}
\newcommand{\vergleichsbilder}{%
\begin{center}\vspace{0.2mm}
\begin{minipage}[c]{0.48\linewidth}\centering
\begin{tikzpicture}[x=1mm,y=1mm,scale=.68,transform shape]
  \path[use as bounding box] (-5,-37) rectangle (45,37);
  % Draufsicht: vollständige Kreisbahn, Kreismitte und tangential ausgerichteter Pkw.
  \draw[Ink,line width=2.2pt] (0,0) circle (31);
  \draw[Ink,line width=2.2pt] (0,0) circle (39);
  \draw[Muted,dashed,line width=.6pt] (0,0) circle (35);
  \fill[Navy] (0,0) circle (1.6);
  \node[anchor=north,font=\scriptsize,text=Muted,fill=white,inner sep=1pt] at (0,-2) {Kreismitte};
  \begin{scope}[shift={(35,0)},rotate=90]
    \draw[Ink,fill=Ink,rounded corners=.8pt] (-9,-7.5) rectangle (-5,7.5);
    \draw[Ink,fill=Ink,rounded corners=.8pt] (5,-7.5) rectangle (9,7.5);
    \draw[Navy,fill=Sky,line width=1pt,rounded corners=2pt] (-14,-6) rectangle (14,6);
    \draw[Blue,fill=white,line width=.7pt,rounded corners=1pt] (4,-4.8) rectangle (10,4.8);
    \draw[Blue,fill=white,line width=.7pt,rounded corners=1pt] (-10,-4.8) rectangle (-5,4.8);
    \draw[Muted,line width=.5pt] (-2,-5)--(-2,5);
    \fill[Gold] (12,-4) rectangle (14,4);
  \end{scope}
  \draw[-{Latex[length=2mm]},Blue,line width=.8pt] (35,-11) -- (35,14);
  \node[font=\scriptsize,text=Blue,anchor=west] at (36,14) {Bewegungsrichtung};
\end{tikzpicture}
\end{minipage}\hfill
\begin{minipage}[c]{0.48\linewidth}\centering
\begin{tikzpicture}[x=1mm,y=1mm,scale=.68,transform shape]
  \path[use as bounding box] (0,-9) rectangle (116,36);
  \draw[Navy,line width=1.1pt] (52,34)--(52,26);
  \draw[Navy,line width=1.3pt] (30,26)--(52,26);
  \draw[Ink,line width=1.4pt] (30,26)--(68,1);
  % Child: head, torso, arms and legs seated on a small swing seat.
  \fill[Gold] (74,12) circle (4.2);
  \draw[Ink,line width=1.8pt,rounded corners=1pt] (72,8)--(69,2);
  \draw[Ink,line width=1.5pt] (71,6)--(78,3);
  \draw[Ink,line width=1.5pt] (69,2)--(79,1) (69,2)--(66,0);
  \draw[Blue,line width=1.2pt] (64,0)--(82,0);
  \draw[Blue,line width=1pt] (66,0)--(66,-3) (80,0)--(80,-3);
\end{tikzpicture}
\end{minipage}\vspace{0.3mm}\end{center}}
\newcommand{\skizzeparallelogramm}{%
\begin{center}\begin{tikzpicture}[x=.56mm,y=.56mm,>=Latex,font=\small]
  \path[use as bounding box] (0,-5) rectangle (125,63);
  % Rotation axis, boom, chain and seated child.
  \draw[Navy,line width=1.1pt] (20,4)--(20,60);
  \draw[Navy,line width=1.1pt] (13,60)--(27,60);
  \draw[Navy,line width=1.3pt] (20,56)--(61,56);
  \coordinate (H) at (61,56); \coordinate (P) at (85,32);
  \draw[Ink,line width=1.25pt] (H)--(P);
  \draw[Blue,line width=1pt,fill=Sky] (82,29) rectangle (98,23);
  \fill[Gold] (94,35) circle (2.5);
  \draw[Ink,line width=1.2pt] (92,33)--(91,29)--(96,27);
  \draw[Ink,line width=1pt] (91,29)--(86,27) (96,27)--(100,25);
  \draw[Blue] (61,56) ++(0,-8) arc[start angle=-90,end angle=-45,radius=8];
  \node[Blue] at (66,47) {$\alpha$};
  % All three forces share P. The dashed edges close the true force parallelogram.
  \draw[-{Latex[length=2.2mm]},Navy,line width=1.1pt] (P)--(85,8) node[right=1pt] {$\FG$};
  \draw[-{Latex[length=2.2mm]},Blue,line width=1.1pt] (P)--(61,56) node[above left=1pt] {$\FS$};
  \draw[-{Latex[length=2.2mm]},Navy,line width=1.2pt] (P)--(61,32) node[midway,below=1pt] {$\FZ$};
  \draw[Muted,densely dashed,line width=.8pt] (61,56)--(61,32) (85,8)--(61,32);
\end{tikzpicture}\end{center}}
\begin{document}
\kopf{Kräfte beim Kettenkarussell \textperiodcentered{} Variante 2}
\frage{1}{Winkelgeschwindigkeit bestimmen}
Auf einen Karussellsitz mit Kind ($m=72\,\mathrm{kg}$) wirkt bei einem Bahnradius von $r=5{,}2\,\mathrm m$ eine Zentripetalkraft von $\FZ=0{,}31\,\mathrm{kN}$. \textbf{Berechne} die Winkelgeschwindigkeit $\omega$.
\antwortfeld{6}{$\FZ=0{,}31\,\mathrm{kN}=3{,}1\cdot10^2\,\mathrm N$. Mit $\FZ=m\omega^2r$ folgt $\omega=\sqrt{\frac{\FZ}{mr}}=\sqrt{\frac{3{,}1\cdot10^2\,\mathrm N}{72\,\mathrm{kg}\cdot5{,}2\,\mathrm m}}=0{,}909940\ldots\,\frac{1}{\mathrm s}\approx0{,}91\,\frac{1}{\mathrm s}$. Die Eingabewerte haben zwei gültige Ziffern; erst das Endergebnis wird gerundet.}
\frage{2}{Auto und Kettenkarussell vergleichen}
\vergleichsbilder
Ein Auto fährt auf ebener Stra\ss{}e im Kreis; ein Sitz mit Kind kreist an schräg gespannter Kette auf gleichbleibender Höhe. \textbf{Erkläre} für beide Fälle, welche realen Kräfte wirken und wie jeweils die zum Kreismittelpunkt gerichtete Zentripetalkraft entsteht.
\antwortfeld{6}{Beim Auto gleichen sich Gewichtskraft nach unten und Normalkraft nach oben aus. Die zum Kreismittelpunkt gerichtete Haftreibungskraft der Stra\ss{}e auf die Reifen liefert die Zentripetalkraft. Bei Sitz und Kind wirken Gewichtskraft nach unten und Seilkraft entlang der schrägen Kette. Die vertikalen Anteile gleichen sich aus; die waagerechte Resultierende zeigt zur Drehachse und ist die Zentripetalkraft. Sie ist keine zusätzliche reale Kraft.}
\fuss{1}
\newpage
\kopf{Auslenkwinkel \textperiodcentered{} Sitzreihe bewerten \textperiodcentered{} Variante 2}
\frage{3}{Auslenkwinkel am Kettenkarussell}
\skizzeparallelogramm
\textbf{Begründe}, ob der Auslenkwinkel $\alpha$ zunimmt, wenn die Winkelgeschwindigkeit $\omega$ zunimmt.\par
{\small\color{Muted}Bilde einen geeigneten Tangens als Ansatz.}\par
\antwortfeld{5}{Im Kräfteparallelogramm gilt $\tan\alpha=\frac{\FZ}{\FG}$. Bei gleicher Masse bleibt $\FG=mg$ gleich. Eine grö\ss{}ere Winkelgeschwindigkeit erfordert eine grö\ss{}ere Zentripetalkraft $\FZ=m\omega^2r$; deshalb wächst $\tan\alpha$ und damit $\alpha$. Der weiter ausgelenkte Sitz hat zugleich einen grö\ss{}eren Bahnradius $r$, der den Kraftbedarf weiter erhöht.}
\frage{4}{Äu\ss{}eren Sitz bewerten}
\begin{minipage}[t]{0.71\linewidth}
  \vspace{0pt}
  Anna möchte Stella überreden, eine Runde im Kettenkarussell mit ihr zu fahren. \textbf{Bewerte} folgende Aussage aus physikalischer Sicht.\par
  \antwortfeld{7}{Annas Aussage ist falsch. Beide Sitze haben dieselbe Winkelgeschwindigkeit $\omega$, weil sie sich mit demselben Karussell drehen. Der äu\ss{}ere Sitz hat den grö\ss{}eren Bahnradius $r$ und wegen $\tan\alpha=\frac{\omega^2r}{g}$ den grö\ss{}eren Auslenkwinkel $\alpha$. Mit $v_{\mathrm B}=\omega r$ ist seine Bahngeschwindigkeit $v_{\mathrm B}$ \textbf{grö\ss{}er}, nicht geringer. Für dieselbe Person ist auch die Zentripetalkraft $\FZ=m\omega^2r$ au\ss{}en grö\ss{}er.}
\end{minipage}\hfill
\begin{minipage}[t]{0.25\linewidth}
  \vspace{1mm}\centering
  \begin{tikzpicture}
    \node[draw=Blue,fill=Sky,rounded corners=2mm,line width=.7pt,
      text width=34mm,align=center,inner sep=2mm,font=\small,text=Ink,execute at begin node={\hyphenpenalty=10000\exhyphenpenalty=10000}] (bubble)
      {\glqq Steig in den äu\ss{}eren Sitz. Dann ist die Geschwindigkeit viel geringer!\grqq};
  \end{tikzpicture}
\end{minipage}\par
\fuss{2}
\end{document}
